


\section{Introduction}

Association schemes are objects of central importance in algebraic combinatorics.
For an introduction to association schemes, the reader could refer to any of~\cite{CP,BI,Z,B,E}.

All our association schemes are symmetric.
If $\X, \Y$ are association schemes on a common vertex set then we write $\X \le \Y$ if $\X$ refines $\Y$ as a partition.
A fusion of an association scheme is a coarsening which is again an association scheme.
{Notice that some authors, for example \cite{KK}, use an opposite order on the set of association schemes. Our choice was motivated by keeping notation consistent with \cite{B,E}. Notice that with this choice of ordering the inclusion $\X \le \Y$ implies a similar inclusion between the automorphism groups $\Aut{\X}\le \Aut{\Y}$.}

We denote the Hamming scheme of order $m^d$ and rank $d+1$ by $\cH(d, m)$
and the Johnson scheme of order $\binom{m}{k}$ and rank $k+1$ by $\cJ(m, k)$.
The special case $\cH(1, m) \cong \cJ(m, 1)$ is the \emph{trivial scheme}, denoted $\cT_m$.
The $d$th tensor power of an association scheme $\X$ is denoted $\X^d$.
The symmetrized $d$th tensor power of the Johnson scheme~$\cJ(m, k)$ is called the \emph{Cameron scheme} and denoted $\cC(m, k, d)$.

In this paper we classify primitive fusions of $\cJ(m, k)^d$ assuming $m$ is sufficiently large in terms of $k$ and $d$.


\begin{theorem}
    \label{main-thm}
    {For any positive integers $k,d$ there exists a constant $m_0(k,d)$ such that
    any primitive fusion $\X$ of $\cJ(m, k)^d$ with $m \geq m_0(k, d)$ belongs, up to permuting coordinates, to one of the following intervals:
    \begin{enumerate}
        \item $\cJ(m, k)^d \leq \X \leq \cC(m, k, d)$,  
        \item $\T_{M^e}^{d/e} \leq \X \leq \cH(d/e, M^e)$ for some integer $e \mid d$ where $M=\binom{m}{k}$.
    \end{enumerate}
    }
\end{theorem}
